Este informe se centra en desarrollar la competencia práctica en el manejo de los controles fundamentales del osciloscopio digital, una herramienta esencial en la ingeniería eléctrica para el análisis de señales.

El objetivo principal de la experiencia fue adquirir la competencia práctica en el uso del osciloscopio digital. Esto implicó progresar desde un conocimiento nulo hasta el dominio de los controles básicos para estabilizar, visualizar y medir correctamente señales eléctricas. Se puso especial énfasis en la comprensión práctica del sistema de disparo (trigger y nivel), los distintos modos de acoplamiento de entrada (CC, CA, GND) y el ajuste de las escalas vertical (ganancia) y horizontal (base de tiempo).

La metodología empleada fue sistemática y experimental, utilizando un Generador de Señales Rigol DG1022 para crear las señales de prueba y un Osciloscopio Digital Rigol DS4014 para su análisis. Se evitó rigurosamente el uso de la función de auto-configuración (``Auto-setting'') para asegurar una comprensión manual de cada parámetro. El desarrollo abarcó el análisis del sistema de disparo, estudiando el efecto del Nivel, la Pendiente y la Fuente; el estudio de los modos de acoplamiento CC, CA, GND y el modo Invertido; el análisis de ganancias y modos temporales, explorando desde la sub-visualización hasta la saturación en modo $Y-t$ y generando figuras de Lissajous en modo $X-Y$; y finalmente, la medición de una señal transitoria configurando el modo de disparo ``Single''.

Se demostró empíricamente que para obtener una traza estable, el nivel de trigger debe estar configurado dentro del rango de amplitud de la señal de referencia. Se comprobó que el acoplamiento CC permite visualizar la señal completa (AC + DC), mientras que el acoplamiento CA filtra la componente DC, centrando la señal en $0 \text{ V}$ pero sin alterar la relación de fase medida. El acoplamiento GND sirvió para establecer la referencia de $0 \text{ V}$. En el análisis de ganancia, se observó cómo ajustes inadecuados llevan a la saturación (``clipping'') o a que la señal sea imperceptible. El modo $X-Y$ permitió visualizar la dependencia directa entre el desfase y la forma de la elipse de Lissajous, así como identificar relaciones de frecuencia. Finalmente, se logró la captura exitosa de una señal transitoria no periódica mediante el uso indispensable del modo de disparo ``Single''.

La comprensión y el ajuste metódico de los sistemas de disparo, acoplamiento y escala son indispensables para garantizar la fiabilidad de las mediciones y evitar interpretaciones erróneas de la señal. El osciloscopio digital se ratifica como una herramienta versátil y esencial, no solo para el análisis de señales periódicas en el dominio del tiempo ($Y-t$) y la fase ($X-Y$), sino también para la captura y análisis de fenómenos transitorios únicos, gracias a su capacidad de almacenamiento digital.